\section{Longest Happy Prefix}
\label{sec:str-longest-happy-prefix}

\problemstatement{A ``happy prefix'' is a non-empty prefix of a string that
is also a suffix (excluding the whole string). Return the longest happy
prefix.}

\begin{patternbox}
The longest happy prefix is precisely the longest proper border of the
string -- which is the \textbf{last value of the KMP prefix function}. No
extra work beyond building the LPS array is needed.
\end{patternbox}

\subsection*{The border is the answer}

The prefix function $\pi[i]$ gives the longest proper border of the prefix
ending at $i$. For the whole string, $\pi[n-1]$ is the length of its longest
proper prefix that is also a suffix -- the definition of the longest happy
prefix. So compute the LPS array and return the prefix of that length.

\begin{ideabox}[{Longest Happy Prefix = Longest Border = $\pi[n-1]$}]
Prefix-that-is-also-suffix is exactly ``border,'' and the prefix function's
final entry stores the longest border of the entire string. The whole
problem is a single LPS computation and a substring.
\end{ideabox}

\subsection*{Algorithm}

\begin{enumerate}
  \item Compute the prefix function $\pi$ of the string.
  \item Return the prefix of length $\pi[n-1]$.
\end{enumerate}

\subsection*{Dry run}

\dryrun{\texttt{"ababab"}.}

\begin{itemize}
  \item $\pi=[0,0,1,2,3,4]$; $\pi[5]=4$.
  \item Longest happy prefix is \texttt{"abab"} (length $4$).
\end{itemize}

\subsection*{C++ implementation}

\begin{lstlisting}
#include <bits/stdc++.h>
using namespace std;

string longestPrefix(const string& s) {
    int n = s.size();
    vector<int> pi(n, 0);
    for (int i = 1; i < n; i++) {
        int j = pi[i - 1];
        while (j > 0 && s[i] != s[j]) j = pi[j - 1];
        if (s[i] == s[j]) j++;
        pi[i] = j;
    }
    return s.substr(0, pi[n - 1]);                 // longest border
}

int main() {
    cout << "[" << longestPrefix("ababab") << "]\n";   // [abab]
    cout << "[" << longestPrefix("leetcode") << "]\n"; // []
    return 0;
}
\end{lstlisting}

\begin{complexitybox}
\textbf{Time Complexity.} $O(n)$. \textbf{Space Complexity.} $O(n)$ for the
prefix-function array.
\end{complexitybox}

\begin{takeawaybox}
The longest happy prefix is just the last value of the KMP prefix function
-- the longest border of the whole string. Recognising ``prefix that is also
a suffix'' as a border collapses the problem to one LPS computation.
\end{takeawaybox}
